% ============================================================
% SECTION 4 — Political Openings & Political Traps (Chapters 4 & 6)
% ============================================================
\section{Political Openings}

\begin{frame}{al-Sadiq: An Opening Is Not a Calling}
\begin{itemize}[itemsep=5pt]
  \item Umayyad collapse creates a real, tempting opening for the Alawid family
  \item Abbasid network organizes under the Prophet's Household's \emph{name} -- without intending to install that Household in power
  \item Abu Salmah sends the \emph{identical} secret offer to both al-Sadiq and 'Abd Allah al-Mahd at once
\end{itemize}
\bigquote{[A poem warning of] the danger of someone doing the work while another reaps the outcome.}{al-Sadiq's reply, after burning Abu Salmah's letter unread}
\delivernote{Ask: why does burning the letter \emph{unread} send a stronger message than reading it and then declining?}
\end{frame}

\begin{frame}{Support the Cause, Reject the Claim}
\begin{itemize}[itemsep=5pt]
  \item Abwa' gathering floats Muhammad Nafs al-Zakiyyah as a possible Mahdi -- on name and coincidence, not doctrine
  \item al-Sadiq refuses the \emph{Mahdi} identification specifically, while saying he would support the movement if genuinely conducted under enjoining good and forbidding evil
  \item Later events (Nafs al-Zakiyyah's death rather than accession) are treated as vindicating al-Sadiq's skepticism
\end{itemize}
\end{frame}

\begin{frame}{al-Rida: Five Competing Theories for Ma'mun's Motive}
\centering\scriptsize
\renewcommand{\arraystretch}{1.5}
\begin{tabular}{|p{3.6cm}|p{6.8cm}|}
\hline
\textbf{Theory} & \textbf{Core claim} \\\hline
Sincere battlefield vow & Vowed to restore the caliphate if granted victory over Amin \\\hline
Fadl ibn Sahl's initiative & Fadl's own project -- al-Rida himself distrusted Fadl \\\hline
Court Iranian sympathy & Appointment courts long-standing Iranian loyalty to the Household \\\hline
Defuse Alawid rebellion & Undercuts rebels' claim the family is shut out of power \\\hline
Neutralize by co-optation & Entangle al-Rida's legitimacy with the government's own failures \\\hline
\end{tabular}
\vspace{0.3cm}
\bigquote{The evidence does not establish any one of these with absolute certainty.}{Mutahhari, explicitly declining to force a verdict}
\end{frame}

\begin{frame}{'Ali ibn Yaqtin's Dilemma, Restated for Ma'mun}
\begin{itemize}[itemsep=5pt]
  \item Summoned via a route deliberately avoiding Shi'ah-populated Kufah
  \item Refuses repeatedly -- accepts only under \emph{direct threat to his life}
  \item Non-interference condition: no judgments, appointments, dismissals, administration
  \item 'Id prayer: crowd's enthusiasm so intense, Ma'mun's own government halts its crown prince mid-procession
\end{itemize}
\bigquote{Either it was rightfully his, in which case he had no right to surrender it -- or it wasn't, in which case he had no right to transfer it.}{al-Rida's dilemma for Ma'mun, on the outright transfer proposal}
\delivernote{The formal acceptance speech pointedly withheld conventional thanks to Ma'mun -- a small, deliberate act of resistance inside an act of outward compliance.}
\end{frame}

\begin{frame}{wilayat al-ja'ir: The Real Test}
\textbf{wilayat al-ja'ir} -- holding office under an unjust ruler -- ordinarily prohibited, but conditional on purpose and effect.
\begin{center}
\Large
\boxed{\text{Whose interest does the cooperation actually serve?}}
\end{center}
\vspace{0.3cm}
\begin{itemize}[itemsep=4pt]
  \item Self-interest condemns it -- Safwan Jammal
  \item Service to the community can justify it -- 'Ali ibn Yaqtin
  \item Coerced survival, stripped of substance -- al-Rida's crown-princeship
\end{itemize}
\end{frame}

\qaframe{Political Openings \& Traps}{%
  \qa{Compare al-Sadiq's withdrawal (Chapter 4) with al-Hassan's peace (Chapter 2). Are they the same kind of decision?}
     {They share a structural logic (decline a fight you'd lose or that costs more than it gains), but al-Hassan negotiates \emph{away from} power he already held; al-Sadiq refuses to even engage with an unproven offer for power he never held. Al-Hassan trades real strength carefully; al-Sadiq declines an offer he judges hollow from the outset.}
  \qa{Item 14 (the non-interference condition) is arguably the single most important fact in the al-Rida chapter. Why does one negotiated clause carry so much interpretive weight?}
     {It converts an ambiguous title into a specific, limited arrangement. Without it, ``al-Rida accepted the crown-princeship'' reads as full cooperation; with it, the identical acceptance reads as a title under duress with every substantive function refused.}
  \qa{State the Safwan/Ibn Yaqtin dividing line in one sentence, then test it against al-Sadiq's refusal in Chapter 4.}
     {Cooperation is judged by whose interest it serves. Al-Sadiq's refusal fits the same logic in reverse: he isn't refusing cooperation with existing power, he's refusing a proposed transfer of power from operators whose sincerity he doubts -- the offer fails the purpose-and-effect test before cooperation could even begin.}
}
